\section{Introduction}

Modern Large Language Models (LLMs) require a complex training pipeline with different algorithms from pre-training to post-training~\citep{team_kimi_2026, deepseek-ai_deepseek-v4_2026, glm-5-team_glm-5_2026}. These algorithms are motivated by different intuitions. For example, the objectives of pre-training and Supervised Fine-Tuning (SFT) arise directly from distribution matching under the cross-entropy loss; On Policy Distillation (OPD) is constructed around on-policy sampling~\citep{lochab_uniform-correct_2026}; Reinforcement Learning (RL) aims to maximize the expected reward~\citep{williams_simple_1992}; and MaxRL can be viewed as arising from maximum-likelihood estimation of correct solutions~\citep{tajwar_maximum_2026}. Our goal is to develop a unified characterization of these seemingly different algorithms in the context of language models. Such a unified perspective may help disentangle the complexity underlying these methods, facilitate the theoretical study of their properties, and provide new insights into their connections and behavior.

Many of the training algorithms of LLMs already admit a distribution matching interpretation~\citep{korbak_reinforcement_2022, khalifa_distributional_2021}, whereas RL is often viewed fundamentally distinct. In this work we first studied the property of an ideal update step in RL and show that, under binary rewards, it preserves the conditional distribution associated with a given reward value. Moreover, if we consider the forward KL divergence between the current model and this conditional distribution, the resulting training direction coincides with that of RL on a fixed sample. Similar results can also be extended to RL with nonbinary rewards.

This provides a unified perspective that places the mainstream training algorithms for LLMs within a common distribution-matching framework. Analyzing how the target distribution changes and affects the learning dynamics under different conditions can yield a variety of interesting theoretical results that characterize the properties of these algorithms.

Depending on whether the target distribution being optimized depends on the current model distribution or not, these algorithms can be further categorized into iterative and non-iterative methods. In what follows, we progressively examine the properties of three classes of algorithms, from simpler to more complex cases: non-iterative methods, iterative methods with zero-mean gradients, and iterative methods with nonzero-mean gradients.

For the first class, we mainly study OPD and SFT-based classical Knowledge Distillation ~\citep{hinton_distilling_2015} algorithms. By analyzing the variance induced by different target distributions under different assumptions, we roughly classified the boundary of effectiveness of these algorithms, providing a theoretical explanation for previous empirical observations regarding the stability of distillation algorithms~\citep{li_rethinking_2026, oh_kl_2026}.

The second class primarily consists of the iterative self-distillation and reversed iterative self-distillation algorithms that we introduce. These methods lie between distillation and RL, while retaining theoretical properties that are more amenable to analysis. Within this class, we identify an intriguing \emph{variance accumulation} phenomenon in iterative training, which leads to a monotonic decrease in the expected entropy of the target distribution. In contrast to previous local results derived from training dynamics~\citep{cui_entropy_2025, ren_learning_2025, jin_revisiting_2026}, our analysis further shows that, although the model entropy may either increase or decrease after an individual update, its long-term behavior exhibits a systematic downward trend under the variance accumulation.

The third class mainly consists of RL and MaxRL. We first provide further empirical evidence that variance accumulation affects entropy in RL-type algorithms, and then analyze the deterministic bias induced by the geometric properties of the model parameterization. Based on these observations, we derive several principled algorithmic corrections, which may also provide useful guidance for the development of future training algorithms. In realistic training settings, however, the target distribution invariance of RL and MaxRL under idealized conditions is no longer preserved. Their target distributions instead exhibit collapse, driven in part by effects such as variance accumulation. Such target distribution drift serves as an important manifestation of the discrepancy between practical training dynamics and their idealized counterparts, offering new insights into the sources of this discrepancy and potential directions for algorithmic improvement.

%In this work we find that all these algorithms can be naturally interpreted as minimizing the divergence between the current model and a target distribution either explicitly or implicitly. Analyzing how the target distribution changes and affects the learning dynamics under different conditions can yield a variety of interesting theoretical results that characterize the properties of these algorithms.

% These algorithms can be further classified into two categories. We call an algorithm iterative if the target distribution depends on the current model distribution, and noniterative otherwise. We primarily studied three different settings from simple to complex.

% \subsection{Noniterative Methods}
%
% We begin by studying distillation methods with a fixed target distribution.
% %Here "noniterative" means that the target distribution is not updated through training steps.
% We find that under the \emph{proximal limit}, where the student model coincides with the target, both forward Knowledge Distillation (KD) and OPD objective has zero mean. We can further prove that OPD has zero variance at the proximal limit while the variance of KD is nonzero. The variance analysis further indicates that the estimation error of OPD is strictly lower than that of KD when the student distribution is close to the teacher. This provides a theoretical guarantee of the effectiveness boundary of these methods, that OPD should be more performant than KD in the proximal case, while being less performant in the distal regime.
%
% \subsection{Iterative Methods with Zero Mean Gradients}
%
% Nonzero gradient-estimation variance at the proximal limit has further consequences for iterative methods, in which the target distribution depends on the current policy. Our analysis shows that ideal updates in distribution space preserve the target distribution in self-distillation and binary-reward RL. In practice, however, gradient-estimation noise can cause the target distribution to drift after each training step. Under local concavity conditions, we show that this \emph{variance accumulation} causes expected entropy to decrease over training, providing a possible mechanism for entropy collapse. Entropy dynamics have been widely studied in RL, with many existing analyses focusing on individual updates~\citep{cui_entropy_2025, jin_revisiting_2026, wang_entropy_2026}. Our analysis extends this perspective to the effects of repeated updates and offers an explanation for long-term behavior.
%
% We design iterative self-distillation and reversed iterative self-distillation experiments to examine the effects of variance accumulation, and observe counter-intuitive results consistent with our theoretical predictions.
%
% \subsection{Iterative Methods with Nonzero Mean Gradients}
%
% RL is a representative iterative distribution-matching method with a nonzero mean gradient. We first show that, when the failure-conditioned distribution has at least as much entropy as the success-conditioned distribution, the target has lower entropy than the current policy. This provides a natural source of entropy reduction during RL training. Such a reduction does not necessarily imply the loss of diversity among correct responses observed in previous work~\citep{yue_does_2025, tajwar_maximum_2026, jin_revisiting_2026}. Our focus is on the deterioration of the target distribution itself, which can arise from both variance accumulation and geometric bias induced by the model parameterization. We conduct experiments to assess the effects of variance on RL algorithms and derive a correction to standard RL through an analysis of this geometric bias. The correction exhibits advantages over standard Group Relative Policy Optimization (GRPO).

% \subsection{Summary of Contributions}

Our contributions are threefold:

\begin{itemize}
\item We introduce a unified theoretical framework that interprets learning algorithms across the LLM training pipeline through distribution matching. By focusing on target distributions and their evolution, this framework clarifies the connections between algorithms and provides a systematic approach to analyzing their properties.
\item We theoretically characterize the regimes in which OPD and SFT have lower gradient-estimation variance, through analysis at the proximal limit and under specified conditions in the distal regime. These results can guide the application of distillation methods and inform future algorithm design.
\item We identify variance accumulation and geometric bias as mechanisms that can drive target-distribution drift and reduce diversity in iterative RL. This analysis motivates a bias correction for RL and suggests directions for future research.
\end{itemize}
